\documentclass[final,5p,times]{elsarticle}
\usepackage[T1]{fontenc}
\usepackage{amsmath,amssymb,booktabs,graphicx,hyperref}
\usepackage{pgfplots}
\usepgfplotslibrary{groupplots}
\pgfplotsset{compat=1.18}
\hypersetup{hidelinks}
\urlstyle{same}
\setlength{\emergencystretch}{2em}
\journal{Array}
\begin{document}
\begin{frontmatter}
\title{Bounded split optimization in decision trees and forests: gains, costs, and limits}
\author[1]{Mathias Valla\corref{cor1}}
\ead{mathias.valla@gmail.com}
\cortext[cor1]{Corresponding author.}
\affiliation[1]{organization={Chaire DIALog, Institut Louis Bachelier},city={Paris},country={France}}
\begin{abstract}
Bounded split optimization is an established response to myopic decision-tree induction, but its predictive gains must be assessed alongside its computational costs. We study single trees, pure bootstrap forests, and forests allocating different search horizons across members. Repeated paired partitions of a fixed classification benchmark compare depth-matched, unpruned learners at shallow, deeper, and unrestricted capacities, with and without feature subsampling. A separate training-only selection experiment applies shared structural grids to conventional and mixed-capable forest families. Sparse replacements improve specific shallow architectures, but the effects change with capacity and feature policy. The inclusive selected family has higher mean accuracy than selected random forests, alongside heterogeneous dataset effects, sensitivity to benchmark weighting, and substantially greater selection and refit costs. Its choices include pure forests, so this gain is not solely evidence for mixing. Mean-effect intervals and adjusted location tests are interpreted separately. The results support conditional benefits on a limited, previously explored sample, not a new induction algorithm, universal accuracy--cost superiority, or a validated per-dataset time-budget policy.
\end{abstract}
\begin{keyword}
Decision trees \sep Bootstrap ensembles \sep Bounded split optimization \sep Computational cost \sep Statistical evaluation
\end{keyword}
\end{frontmatter}

\section{Introduction}
Decision trees choose a sequence of splits, but most induction methods optimize each split separately. CART minimizes the weighted impurity of the two immediate children and then recurses \cite{breiman1984cart}. A candidate with little immediate gain may enable a useful interaction one generation later. Bounded descendant optimization addresses this mismatch by scoring the first split through the best local subtree it enables.

This idea is established. The question here is empirical: \emph{does bounded split optimization improve held-out accuracy enough to justify its measured fitting cost, and does applying it to a minority of ensemble members help?} We use $k$-sighted to denote an integer local horizon of $k$ split generations, including the current split. A mixed forest has members with different integer horizons; its average horizon is descriptive rather than a fractional optimization rule.

We examine two distinct questions. Under a fixed architecture, does increasing the horizon of all members or replacing only a minority improve prediction? After training-only tuning, does a mixed-capable selection family improve on a conventional forest family given the same common hyperparameter choices? Repeated paired partitions separate these comparisons, with matching depth limits and no pruning. Shallow, deeper, and unrestricted capacities test whether an apparent search benefit depends on constraining the trees. Measured selection work and final fitting cost are reported separately. Aggregate accuracy--time plots do not validate a model chosen under a per-dataset time budget. The study concerns a fixed sample of small PMLB tasks, not general machine-learning workloads.

The shallow regime is deliberate, not a proposed default for random forests. Limiting member size can regularize an ensemble, with conditional predictive benefits documented particularly in noisy regression settings \cite{zhou2023trees,duroux2018impact}. At depth three each binary member has at most eight leaves; that bound does not limit the whole forest to eight prediction regions. We retain this capacity-limited question while testing deeper alternatives rather than assuming that shallow members are universally sufficient.

The contribution is an evaluation of computation allocation, not a new bounded-search algorithm or an optimality certificate. The implementation and reproducibility guide are available at \href{https://github.com/MathiasValla/non-greedy-decision-trees}{the project repository}. The one-sighted baseline in the new protocol is library CART, rather than an independently implemented greedy learner with different stopping behaviour.

The completed evaluation finds gains for specific shallow replacements and a positive mean advantage for an expanded training-selected family. Neither result establishes a general accuracy--cost improvement: effects vary across tasks and settings, selected models are not always mixed, and the expanded selection spaces require substantially more computation. These distinctions guide the analysis rather than treating every positive mean as evidence for a practical replacement of greedy forests.

\section{Related work}
Greedy induction and bootstrap aggregation are standard foundations \cite{breiman1984cart,quinlan1986induction,breiman1996bagging}. Conventional random forests additionally randomize the features available at each split \cite{breiman2001random}; a bootstrap ensemble using all features is not that conventional baseline.

Non-greedy induction predates the present study. Norton \cite{norton1989generating} considered downstream tree quality, while Ragavan and Rendell \cite{ragavan1993lookahead} used lookahead for feature construction. Murthy and Salzberg's LOOK procedure directly evaluates a split with optimally split children, corresponding to the two-generation case here \cite{murthy1995lookahead}. Their findings already question whether added computation reliably improves prediction. Later work revisited pathology and anytime computation allocation \cite{elomaa2005lookahead,esmeir2007anytime}. Norouzi et al. \cite{norouzi2015efficient} jointly optimize oblique split and leaf parameters, a different form of non-greediness.

The closest forest precedent is the stepwise lookahead forest of Donick and Lera \cite{donick2021stepwise}. It optimizes two-level blocks and studies interactions in noisy data. Our builder instead commits one root split and re-solves at each actual child; our ensemble question concerns allocating different horizons across members. We do not claim priority for lookahead forests or heterogeneous ensembles. Randomized splits \cite{geurts2006extremely} and random final depths \cite{horn2026randomdepth} are other established ensemble design choices. Heterogeneity alone is not evidence of a beneficial diversity mechanism \cite{kuncheva2003diversity}.

Modern optimal-tree methods exploit mathematical programming, bounds, caching, and separable subproblems \cite{bertsimas2017optimal,verwer2019learning,hu2019optimal,lin2020generalized,aglin2020caching,demirovic2022murtree,vanderlinden2023separable,aghaei2025strong}. Reference ensembles can guide restricted search \cite{mctavish2022reference}. SPLIT and LicketySPLIT combine non-greedy choices with greedy continuation and distinguish search depth from final depth \cite{babbar2025nearoptimal}. Their objectives and continuation rules differ from truncated terminal impurity. Recent comparisons also emphasize that objectives, size constraints, and tuning affect conclusions about greedy versus optimal trees \cite{vanderlinden2026optimalgreedy}. We do not benchmark these solvers here and cannot infer that our measured cost is unavoidable for every non-greedy method.

\subsection{Tree-size regularization}
Shallow forests have a rationale beyond making search tractable, but the evidence is conditional. Zhou and Mentch \cite{zhou2023trees} study tree size as random-forest regularization and report benefits for smaller trees at low signal-to-noise ratios. Duroux and Scornet \cite{duroux2018impact} compare subsampling and early termination theoretically for quantile regression forests and empirically for Breiman-style regression forests. Their leaf-count and resampling settings differ from our bootstrap depth-three Gini classifiers; neither study identifies depth three as a general optimum. Competitive shallow single-tree classification results \cite{bertsimas2017optimal} motivate a separate capacity-constrained question, not a guarantee for forests. Consistency theory with controlled but growing leaf counts \cite{scornet2015consistency} likewise does not establish consistency for fixed depth three. We therefore situate the shallow experiment within tree-size regularization research without substituting it for a comparison with deep forests. Results about shallow learners in sequential boosting are not evidence for independently fitted bootstrap ensembles.

\section{Method}
\subsection{Horizon and final depth}
Let $S$ be the examples reaching a node, $I(S)$ its Gini impurity, and $W(S)$ its total weight. Let $D$ be the maximum number of splits along a final root-to-leaf path and $d$ the current depth. At that node the actual search horizon is $h=\min(k,D-d)$; when $D$ is unrestricted, $h=k$. Increasing $k$ changes search, whereas increasing $D$ changes capacity.

Define the optimal local terminal score recursively:
\begin{align}
 C_0(S)&=W(S)I(S),\\
 C_h(S)&=\min\left\{C_0(S),\min_{s\in\mathcal S(S)}
 \big[C_{h-1}(S_l(s))+C_{h-1}(S_r(s))\big]\right\}.
 \label{eq:recursion}
\end{align}
Here $\mathcal S(S)$ contains admissible axis-aligned midpoint thresholds between distinct observed feature values, subject to the leaf, depth, and feature constraints. The selected root minimizes the second term when it is no worse than stopping, up to numerical tolerance. Feasible zero-gain splits are allowed, as in unpruned CART. Only that root is committed. The configured horizon is then restored at each actual child. This is receding-horizon induction, not a single depth-$k$ fit followed by an unchanged completion.

For a fixed node, expanding the feasible horizon cannot worsen the optimal terminal \emph{training impurity}. This does not guarantee better test accuracy or monotone performance of the final receding-horizon trees. Gini is not misclassification error or a sparsity-penalized loss. With all features and $k\geq D$ at the root, exhaustive search solves the depth-bounded terminal-impurity problem, subject to the stated constraints and numerical tolerance. With feature subsampling, optimization is conditional on the seeded, node-dependent feature policy, not over every axis-aligned subtree. Thus $k=D=3$ with all features is a special shallow full-horizon case, not evidence of scalable global optimization.

\subsection{Implementation and stopping}
The isolated Cython scorer evaluates uncapped numerical thresholds and reuses sorted prefix class counts for the terminal split calculation. It does not introduce branch-and-bound or a global optimality certificate. Pure or infeasible nodes stop, but a feasible split is not rejected solely for having zero gain. Both implementations receive float32 predictors, Gini impurity, the same permitted total depth, split and leaf minima, feature setting, and seed. The minimum impurity decrease and minimum leaf-weight fraction are zero. CART uses \texttt{ccp\_alpha=0}; the sighted builder has no post-fit pruning. With $p$ predictors, square-root subsampling initially considers $\max(1,\lfloor\sqrt p\rfloor)$ features at a node. This is not a strict inspection cap: additional features are examined if the initial randomized subset offers no admissible partition, including when its features are constant. Threshold rounding, tie-breaking, and random-feature ordering can still differ. Horizon is the deliberate algorithmic change, not a change in capacity or pruning.

The one-sighted constituent is scikit-learn CART. Forest tree seeds and bootstrap multiplicities follow scikit-learn's random-forest construction. Each member fits the full training matrix using multiplicities as sample weights; zero-weight rows are excluded from the sighted builder's sample-count constraints. This avoids the different minimum-leaf semantics of physically duplicated rows. A retained verification harness checks prepared CART-bank predictions against the actual library random forest across the evaluated ensemble sizes, depths, feature settings, and leaf minima, including constant predictors and bags missing a rare class, with absolute tolerance $10^{-12}$ and zero relative tolerance. This controlled check does not establish identity between CART and the farther-sighted implementations.

Each mixed-forest bootstrap position contains exactly one tree. A farther-sighted member replaces CART at that position; it does not add a second tree trained on an already-counted draw. The paired fits use the same weighted draw and seed. Label-independent random permutations of successive 20-position blocks assign disjoint horizon groups. This gives exact proportions and nested prefixes at every evaluated ensemble size. Class probabilities are aligned to the common dataset class ordering before averaging.

The shorthand $\bar k=\sum_j p_j k_j$ describes a mixture. It is not a sufficient description of its search or cost: $90\%$ one-sighted plus $10\%$ two-sighted and $95\%$ one-sighted plus $5\%$ three-sighted both have $\bar k=1.1$. We specify proportions and ensemble size directly.

\subsection{Search cost}
Counting complete split assignments is not the runtime of Eq.~\eqref{eq:recursion}. Additivity allows the child optima to be computed separately. If every recursive call has at most $m$ candidates and $E_h$ counts candidate evaluations, a uniform worst-case abstraction gives
\begin{equation}
 E_0=0,\qquad E_h\leq m+2m E_{h-1},\qquad
 E_h\leq m\frac{(2m)^h-1}{2m-1}.
\end{equation}
The bound is exponential in horizon, before accounting for sorting, partitioning, changing sample sizes, and early stopping. Final fit cost additionally sums searches over committed nodes. A fixed short horizon can therefore be used in deep trees; whether its cost is acceptable is an empirical question. Our timings characterize this implementation and candidate space, not an inherent lower bound for optimized non-greedy induction.

\section{Experimental protocol}
\subsection{Dataset sample and repeated partitions}
We freeze the 57-dataset sample previously used for the mixed-forest grid. This is a convenience sample of tractable PMLB classification tasks \cite{olson2017pmlb}, not a random sample of all classification problems or a fresh external validation collection. Every dataset is retained in the new protocol. Dataset names, raw and used dimensions, class counts, and split hashes are retained. PMLB's supplied numeric representation and categorical encodings are shared by all methods. Loading explicitly applies the same complete-case rule to every family; the manifest records missingness, removed-row counts, and a retained source-index hash. This removes eleven rows from \texttt{penguins}. No imputation, learned encoding, scaling, or candidate cap is introduced.

Each dataset receives five paired stratified $75/25$ outer partitions, with seeds $1000$--$1004$. Estimator and bootstrap seeds vary across repetitions. All methods use the same outer indices. This increases assessment of split stability; overlapping partitions are not independent datasets. The study is an extension of an exploratory benchmark, not preregistered independent confirmation.

\subsection{Models and training-only selection}
Fixed-architecture comparisons use $D\in\{3,6,\infty\}$, minimum leaf size one, and either all features or square-root feature subsampling. Single trees use $k=1,2,3$ with identical common settings and no pruning. Every forest composition is evaluated at $T\in\{20,40,60,100,200\}$. In addition to the three pure horizons, each pair $(1,2)$, $(1,3)$, and $(2,3)$ uses $5\%,10\%,25\%,$ or $50\%$ of the farther-sighted horizon. A sparse three-way composition uses $85\%,10\%,5\%$ of horizons one, two, and three. Thus sixteen explicitly specified compositions span every ensemble size; equal average horizons do not identify identical mixtures.

The tuned experiment is separate. The library RF family, the one-/two-sighted family, and the inclusive one-/two-/three-sighted family all receive training-only inner cross-validation. They share depths $\{3,6,\infty\}$, leaf minima $\{1,5\}$, feature settings $\{\sqrt p,p\}$, and ensemble sizes $\{20,40,60,100\}$. Here RF denotes the library estimator family: it includes both feature-subsampled forests and all-feature bagged CART, not exclusively conventional square-root-feature forests. The 200-member fixed curves are descriptive and not extra tuning candidates for only one family. The RF space has 48 candidates, the one-/two-sighted space 288, and the inclusive space 768; both mixed spaces contain the full RF space. Equal common search grids do not mean equal computational effort, and the larger mixed spaces can incur greater selection uncertainty and cost.

Inner folds are shared across families. Shuffled stratified cross-validation uses $\min(3,c_{\min})$ folds when each observed outer-training class has at least two examples. A singleton or single-class training partition uses ordinary shuffled three-fold cross-validation; this fallback is documented and cannot provide reliable rare-class validation. Shared horizon-specific banks are fitted once per fold and structural configuration, and scored at every admissible size and composition. Highest mean inner accuracy selects, followed by lowest fold-mean fitting work of the candidate's included members and then a stable candidate identifier to break ties. This training-only timing tie-break is hardware-dependent; it is not a resource constraint or the full cost of selection. All three configurations are locked before separate refits and outer-test evaluation \cite{cawley2010selection}. The selected conventional baseline refits the actual library RF. Outer test outcomes never choose candidates.

Three-sighted members are included in both repeated fixed and tuned experiments, not treated only as historical evidence. No fit timeout is imposed. Incremental bank checkpoints permit recovery without deleting unfinished tasks. Completeness requires every planned dataset, repetition, structural block, and selected refit; an error stops execution rather than silently shrinking the sample.

The full analysis checks completeness, reconstructs fixed and inner-validation forest probabilities and constituent fitting costs from retained banks, and independently reconstructs the validation winners. Single-tree and direct selected-refit scores receive range and checkpoint consistency checks, not independent prediction reconstruction. Retained manifests and the analysis certificate identify the protocol, partitions, fit sources, and certified numerical exports; provenance checks do not constitute an independent experimental replication.

\subsection{Timing and outcomes}
Accuracy is the primary metric and balanced accuracy is an imbalance sensitivity. Balanced accuracy follows the library definition: mean recall over classes present in the scoring partition. If a predicted class is absent from that partition, its recall is undefined and excluded from this mean; the corresponding warning is not a fit failure. This convention cannot evaluate recall for absent rare classes. Scores are first averaged over the five repetitions within each dataset, then averaged equally across datasets. We retain all repeat-level scores and selected baseline settings.

Fixed-curve fitting work sums measured \texttt{fit} durations for exactly the included members. Constructor/bootstrap preparation and prediction are measured separately, not silently included in this axis. Fixed single-tree timings also measure \texttt{fit} alone, excluding construction and prediction. Direct selected-model refit wall time includes construction and bootstrap preparation but excludes prediction. We report validation selection work separately from final refit cost. Each tuned family is charged the banks and schedules required by its complete candidate space, including unused members of computed banks. Shared computation is charged fully to each standalone family but counted only once when reporting the physical benchmark work. A sum of recorded stage wall times excludes checkpoint I/O, process startup, and idle gaps; it is not a full uninterrupted end-to-end stopwatch. These experiments do not enforce an equal total computation budget.

Five concurrent shards each use one thread for fitting and numerical libraries on a MacBook Air with an Apple M3 processor (eight CPU cores) and 16 GB memory. Wall times reflect the recorded software and shared hardware load, not isolated-machine speed measurements. Cross-implementation timing comparisons include implementation effects. Close time matches must therefore not be interpreted as guaranteed budget feasibility. We do not optimize a utility weight or claim a validated fixed-time frontier.

OpenAI Codex (developer: OpenAI) assisted development and revision of the benchmark, analysis, and visualization code. Its tool version and model identifier were not recorded in the retained research provenance and are unavailable from those records. Numerical figures derive directly from the retained results through reproducible NumPy, SciPy, pandas, and Matplotlib workflows, with equivalent native PGFPlots plots embedded in this source; no generative image model supplies data or artwork. Exact Python and numerical-library versions are recorded in the run manifests and full-analysis certificate.

\subsection{Uncertainty and multiple comparisons}
For each contrast we average the paired repetition differences within each dataset. A percentile bootstrap of these dataset means, with 20,000 resamples, gives an unadjusted $95\%$ interval for the equal-dataset mean effect. It describes variation across the benchmark tasks under a dataset-exchangeability approximation; it is not a distribution-free confidence interval for all machine-learning applications.

Two-sided Wilcoxon signed-rank tests use dataset-level differences, omit zero differences, and assign $p=1$ for an all-zero contrast \cite{demsar2006comparisons}. They test a signed-rank location hypothesis, requiring the usual symmetry qualification, rather than the same mean estimated by the bootstrap. Holm adjustment \cite{holm1979sequential} covers one prespecified family of eleven contrasts: eight fixed 100-member comparisons of $10\%$ two- or three-sighted replacements with matched RF, at depth three or unrestricted depth and either feature setting; and the three pairwise comparisons of the tuned families. Other depths, sizes, proportions, and single-tree contrasts are descriptive. Intervals are not simultaneous. This contrast family was frozen before the new run; the underlying cohort had already been explored.

Known related GAMETES, MONK, LED, parity, and Cy Young tasks are also averaged within families for a sensitivity analysis with equal family weighting. A second sensitivity removes the explicitly named synthetic/constructed tasks in the released manifest. Both change the weighting or sample; neither selects favourable results. Dataset dependence beyond those known groups remains possible. Failure to detect a difference is not evidence of equivalence.

\section{Results}
% BEGIN GENERATED results_text
% First complete-results writing pass, for independent scientific review.
% Source: paper/array_revision_fair/analysis/full; no new inference or fits.
% Replace the Results body, not the Methods. Table fragments are complete tables.
% Place the existing validated Figure 2 at the insertion marker below.
% Figure 3 below references the existing full-analysis PDF; do not regenerate it.

\subsection{Complete evaluation}
Both stages completed all 57 datasets and five paired outer partitions,
with no missing or skipped task. The fixed stage contains 1,710
depth--feature blocks, 5,130 single-tree scores, and 136,800 forest
scores. The selected stage adds 855 separately refitted forests, three
families for each of the 285 outer tasks. Table~\ref{tab:performance}
reports representative fixed architectures and all selected families.
Differences below are percentage points of accuracy. Intervals are
pointwise, unadjusted $95\%$ paired dataset-bootstrap intervals for
the mean effect. Table~\ref{tab:comparisons} gives the eleven primary
contrasts with one joint Holm adjustment; other comparisons are descriptive.
Scores and costs first average repetitions within datasets, then weight
datasets equally.

% BEGIN GENERATED table_performance
% COMPLETE table environment: replace the old table, not only its row markers.
% 24 fixed rows (all features, leaf minimum one) and three selected-family rows.
% Source row IDs and timing fields are documented in full_interpretation_notes.md.
\begin{table*}[t]
\centering\small
\caption{Performance on 57 datasets and five paired repetitions. Scores,
costs, and realized depths first average within datasets, then equally
across datasets. Fixed rows use all features and minimum leaf size one;
all fixed forests have 100 members. Replacement forests contain 90\%
CART and 10\% of the indicated horizon. Fixed-tree fit measures are
fit-only wall times; fixed-forest measures are selected-constituent fit
sums, both excluding construction/bootstrap preparation and prediction.
Selected-family rows choose depth, leaf size, features (all or square-root),
ensemble size, and admissible composition using training-only validation.
Their fit measure is direct refit wall time including construction/bootstrap
preparation but excluding test prediction and selection. The differing
timing scopes are not interchangeable or total learning costs. Accuracy
and balanced accuracy are proportions.}
\label{tab:performance}
\begin{tabular}{@{}lrrrrr@{}}\toprule
Model & $D$ & Accuracy & Balanced accuracy & Fit measure (s) & Realized depth\\
\midrule
CART tree & 3 & 0.6782 & 0.6241 & 0.0010 & 3.0 \\
Tree $k=2$ & 3 & 0.7059 & 0.6543 & 0.0282 & 3.0 \\
Tree $k=3$ & 3 & 0.7136 & 0.6647 & 1.6845 & 3.0 \\
Pure $k=1$ forest & 3 & 0.7288 & 0.6721 & 0.0677 & 2.9 \\
Pure $k=2$ forest & 3 & 0.7478 & 0.6931 & 1.6807 & 3.0 \\
Pure $k=3$ forest & 3 & 0.7522 & 0.6959 & 80.2307 & 3.0 \\
10\% $k=2$ replacement & 3 & 0.7393 & 0.6835 & 0.2293 & 2.9 \\
10\% $k=3$ replacement & 3 & 0.7415 & 0.6853 & 8.0933 & 2.9 \\
\addlinespace
CART tree & 6 & 0.7282 & 0.6947 & 0.0011 & 5.6 \\
Tree $k=2$ & 6 & 0.7321 & 0.6977 & 0.0569 & 5.9 \\
Tree $k=3$ & 6 & 0.7341 & 0.7009 & 3.7724 & 5.9 \\
Pure $k=1$ forest & 6 & 0.7665 & 0.7297 & 0.0797 & 5.4 \\
Pure $k=2$ forest & 6 & 0.7606 & 0.7201 & 3.7387 & 5.9 \\
Pure $k=3$ forest & 6 & 0.7556 & 0.7166 & 200.3026 & 5.9 \\
10\% $k=2$ replacement & 6 & 0.7692 & 0.7321 & 0.4463 & 5.4 \\
10\% $k=3$ replacement & 6 & 0.7723 & 0.7350 & 20.1196 & 5.4 \\
\addlinespace
CART tree & $\infty$ & 0.7451 & 0.7194 & 0.0017 & 10.1 \\
Tree $k=2$ & $\infty$ & 0.6938 & 0.6578 & 0.1592 & 15.1 \\
Tree $k=3$ & $\infty$ & 0.6622 & 0.6208 & 7.9657 & 17.0 \\
Pure $k=1$ forest & $\infty$ & 0.7752 & 0.7450 & 0.1466 & 8.9 \\
Pure $k=2$ forest & $\infty$ & 0.7069 & 0.6615 & 9.7728 & 13.2 \\
Pure $k=3$ forest & $\infty$ & 0.6784 & 0.6302 & 387.2398 & 14.9 \\
10\% $k=2$ replacement & $\infty$ & 0.7758 & 0.7442 & 1.1200 & 9.4 \\
10\% $k=3$ replacement & $\infty$ & 0.7750 & 0.7430 & 38.5187 & 9.5 \\
\midrule
Selected RF family & selected & 0.7899 & 0.7533 & 0.0626 & 6.5 \\
Selected $k\leq2$ family & selected & 0.7996 & 0.7642 & 0.9518 & 5.5 \\
Selected inclusive family & selected & 0.8056 & 0.7693 & 12.9374 & 5.4 \\
\bottomrule
\end{tabular}
\end{table*}
% END GENERATED table_performance

% BEGIN GENERATED table_comparisons
% COMPLETE table environment: replace the old table, not only its row markers.
% Exactly the eleven rows of full/paired_primary_comparisons.csv, in file order.
\begin{table*}[t]
\centering\small
\caption{All eleven prespecified primary paired comparisons. Fixed rows
replace ten of 100 CART members with the indicated horizon, at minimum
leaf size one and the stated feature setting. Effects and pointwise,
unadjusted $95\%$ mean-effect intervals are in percentage points.
$p_H$ is the two-sided Wilcoxon signed-rank value Holm-adjusted jointly
over these eleven contrasts, not a test of the same mean as the interval.
W/T/L counts dataset-mean wins/ties/losses for the first model. Cost ratios
divide equal-dataset mean costs: selected-constituent fitting work for
the eight fixed rows, direct selected-model refit wall time for the three
selected-family rows. They exclude selection work and have different
timing boundaries. Extra precision retains the near-zero interval endpoint
in the unrestricted square-root-feature horizon-three row.}
\label{tab:comparisons}
\begin{tabular}{@{}lrlrrr@{}}\toprule
Comparison & $D$ & Mean difference [95\% interval] & $p_H$ & W/T/L & Cost ratio\\
\midrule
10\% $k=2$ vs. CART (all) & 3 & $1.06$ [$0.45$, $1.77$] & 0.00996 & 27/21/9 & 3.39 \\
10\% $k=3$ vs. CART (all) & 3 & $1.27$ [$0.61$, $2.02$] & 0.02066 & 28/14/15 & 119.52 \\
10\% $k=2$ vs. CART (sqrt) & 3 & $0.42$ [$0.07$, $0.78$] & 0.09539 & 27/22/8 & 4.84 \\
10\% $k=3$ vs. CART (sqrt) & 3 & $1.04$ [$0.46$, $1.74$] & 0.00150 & 32/19/6 & 106.73 \\
10\% $k=2$ vs. CART (all) & $\infty$ & $0.05$ [$-0.27$, $0.44$] & 1.00000 & 19/20/18 & 7.64 \\
10\% $k=3$ vs. CART (all) & $\infty$ & $-0.02$ [$-0.49$, $0.55$] & 1.00000 & 16/20/21 & 262.81 \\
10\% $k=2$ vs. CART (sqrt) & $\infty$ & $0.22$ [$-0.02$, $0.47$] & 0.75287 & 23/18/16 & 18.90 \\
10\% $k=3$ vs. CART (sqrt) & $\infty$ & $0.28$ [$0.00047$, $0.59$] & 0.34732 & 25/16/16 & 367.04 \\
\midrule
Selected $k\leq2$ vs. RF & selected & $0.97$ [$-0.07$, $2.30$] & 1.00000 & 29/8/20 & 15.20 \\
Selected inclusive vs. RF & selected & $1.57$ [$0.30$, $3.28$] & 0.48838 & 28/11/18 & 206.61 \\
Selected inclusive vs. $k\leq2$ & selected & $0.60$ [$-0.11$, $1.63$] & 1.00000 & 18/23/16 & 13.59 \\
\bottomrule
\end{tabular}
\end{table*}
% END GENERATED table_comparisons

\subsection{All-feature single trees and pure forests}
At $D=3$, mean accuracies for single trees with $k=1,2,3$ are
$67.82\%,70.59\%,71.36\%$. The paired horizon-two/three gains over
CART are $2.77$ points [$0.68$, $5.22$] and $3.54$ points [$1.45$,
$6.01$], but medians are only $0.13$ and $0.40$ points, with
$29/3/25$ and $31/3/23$ wins/ties/losses. Mean fit-only times are
$0.00103$, $0.0282$, and $1.684$ s, respectively. At $D=6$, mean
differences shrink to $0.39/0.59$ points and medians become
$-0.48/-1.27$ points. Without a depth limit, mean differences are
$-5.13$ [$-8.69$, $-1.84$] and $-8.29$ [$-12.45$, $-4.39$]
points. Mean realized depths then increase from $10.1$ for CART to
$15.1/17.0$ for horizons two/three. Larger realized trees and poorer
test scores coincide here, without identifying their cause.

For pure $T=100$ forests with all features, horizon-two/three mean
differences from the matched CART forest are $1.90/2.34$ points at
$D=3$, $-0.58/-1.09$ at $D=6$, and $-6.84/-9.68$ without a depth
limit. Shallow mean fitting work is $0.0677$, $1.6807$, and $80.2307$
s for the three pure horizons. With square-root feature subsampling,
the corresponding mean differences are $2.38/2.85$, $1.27/1.96$,
and $1.08/0.78$ points. The last two intervals are [$0.35$, $1.88$]
and [$-0.75$, $2.44$], with medians $0.29$ and zero points.
These are contrasts with each feature policy's own CART control,
not a formal interaction test or evidence that subsampling universally
improves absolute accuracy. The unrestricted all-feature losses must
not be extrapolated to the square-root-feature setting.

\begin{figure*}[t]
\centering% BEGIN GENERATED figure1
\definecolor{fairC0}{HTML}{222222}
\definecolor{fairC1}{HTML}{1F77B4}
\definecolor{fairC2}{HTML}{D62728}
\definecolor{fairC3}{HTML}{2B8C55}
\definecolor{fairC4}{HTML}{66A61E}
\definecolor{fairC5}{HTML}{C59A00}
\definecolor{fairC6}{HTML}{8C6D31}
\definecolor{fairC7}{HTML}{7B3294}
\definecolor{fairC8}{HTML}{A64C9C}
\definecolor{fairC9}{HTML}{DC74AB}
\definecolor{fairC10}{HTML}{B35806}
\definecolor{fairC11}{HTML}{008B8B}
\definecolor{fairC12}{HTML}{17A8BD}
\definecolor{fairC13}{HTML}{527A9C}
\definecolor{fairC14}{HTML}{7F7F7F}
\definecolor{fairC15}{HTML}{9D8F11}
\begin{tikzpicture}
\begin{groupplot}[group style={group size=2 by 1,horizontal sep=1.1cm},
width=0.44\textwidth,height=5.2cm,grid=major,grid style={gray!15},
tick label style={font=\scriptsize},label style={font=\small},title style={font=\small},
xtick={1,2,3},xlabel={Sight horizon $k$}]
\nextgroupplot[title={(a) Paired single-tree differences},ylabel={Accuracy difference (pp)},
legend to name=fairDepthLegend,legend columns=3,legend style={font=\small,draw=none}]
\addplot+[mark=*,color=fairC1,error bars/.cd,y dir=both,y explicit] table[x=x,y=y,y error minus=low,y error plus=high] {
x y low high
0.940 0 0 0
1.940 2.7690655 2.0892691 2.4517777
2.940 3.5405952 2.092686 2.4678195
};
\addlegendentry{Depth 3}
\addplot+[mark=*,color=fairC3,error bars/.cd,y dir=both,y explicit] table[x=x,y=y,y error minus=low,y error plus=high] {
x y low high
1.000 0 0 0
2.000 0.38921788 2.2884109 2.5959166
3.000 0.58875764 2.504071 2.9363957
};
\addlegendentry{Depth 6}
\addplot+[mark=*,color=fairC2,error bars/.cd,y dir=both,y explicit] table[x=x,y=y,y error minus=low,y error plus=high] {
x y low high
1.060 0 0 0
2.060 -5.1289459 3.5641001 3.2868148
3.060 -8.2912463 4.1597501 3.902189
};
\addlegendentry{Unlimited depth}
\addplot[black!40,dashed,forget plot] coordinates {(0.8,0) (3.2,0)};
\nextgroupplot[title={(b) Single-tree fitting cost},ymode=log,ylabel={Mean fitting time (s)}]
\addplot+[mark=*,color=fairC1] coordinates {(1,0.00102871712) (2,0.0282174026) (3,1.68448913)};
\addplot+[mark=*,color=fairC3] coordinates {(1,0.00112826272) (2,0.0569112694) (3,3.7723653)};
\addplot+[mark=*,color=fairC2] coordinates {(1,0.00169035613) (2,0.159224192) (3,7.96574853)};
\end{groupplot}
\end{tikzpicture}
\par\smallskip\pgfplotslegendfromname{fairDepthLegend}
% END GENERATED figure1
\caption{Depth-matched unpruned single trees with all features. Left: mean paired accuracy differences from CART, with pointwise, unadjusted $95\%$ dataset-bootstrap intervals; these single-tree effects are descriptive, not members of the eleven-test primary family. Right: measured mean single-tree fit duration, excluding construction and prediction. Sight horizon is varied independently of permitted final depth.}
\label{fig:depth}
\end{figure*}

\subsection{Sparse horizon replacements}
At $T=100$, $D=3$, with all features, replacing ten CART members with
horizon-two or horizon-three members gives mean gains of $1.06$ points
[$0.45$, $1.77$] and $1.27$ points [$0.61$, $2.02$]. Both medians
are zero; wins/ties/losses are $27/21/9$ and $28/14/15$.
Their joint-Holm signed-rank values are $p_H=0.00996$ and $0.02066$.
Balanced-accuracy gains are $1.15/1.33$ points. Mean selected-slot
fitting work rises from $0.0677$ s for CART to $0.2293/8.0933$ s,
or $3.39/119.52$ times the control. These are architecture-matched,
not equal-computation gains. The descriptive $85/10/5\%$ composition
has a mean gain of $1.36$ points [$0.52$, $2.30$] at $4.2260$ s.

With square-root features, the two shallow replacement gains are
$0.42$ [$0.07$, $0.78$] and $1.04$ [$0.46$, $1.74$] points,
with $p_H=0.09539$ and $0.00150$. Their medians are zero and
$0.13$ points. Without a depth limit, all four primary replacement
means are small: $0.05/-0.02$ points with all features and
$0.22/0.28$ points with square-root features. All four medians are
zero; adjusted values range from $0.34732$ to $1$. This is not
equivalence evidence. At the descriptive depth-six all-feature
setting, the replacement mean differences are $0.27/0.58$ points,
again with zero medians.

For the two shallow all-feature anchors, equal weighting of the 48
known dataset families reduces the means to $0.77/0.91$ points,
with intervals [$0.25$, $1.45$] and [$0.32$, $1.63$]. On the 38
tasks outside the named synthetic/constructed set, the means are
$0.53/0.63$ points, with intervals [$0.06$, $1.22$] and [$0.07$,
$1.35$]. These sensitivities concern these anchors and do not make
the previously explored convenience sample independent or representative.

\subsection{Descriptive mean-cost curves}
\begin{figure*}[t]
\centering% BEGIN GENERATED figure2
\definecolor{fairC0}{HTML}{222222}
\definecolor{fairC1}{HTML}{1F77B4}
\definecolor{fairC2}{HTML}{D62728}
\definecolor{fairC3}{HTML}{2B8C55}
\definecolor{fairC4}{HTML}{66A61E}
\definecolor{fairC5}{HTML}{C59A00}
\definecolor{fairC6}{HTML}{8C6D31}
\definecolor{fairC7}{HTML}{7B3294}
\definecolor{fairC8}{HTML}{A64C9C}
\definecolor{fairC9}{HTML}{DC74AB}
\definecolor{fairC10}{HTML}{B35806}
\definecolor{fairC11}{HTML}{008B8B}
\definecolor{fairC12}{HTML}{17A8BD}
\definecolor{fairC13}{HTML}{527A9C}
\definecolor{fairC14}{HTML}{7F7F7F}
\definecolor{fairC15}{HTML}{9D8F11}
\begin{tikzpicture}
\begin{groupplot}[group style={group size=2 by 1,horizontal sep=0.95cm},
width=0.445\textwidth,height=5.5cm,grid=major,grid style={gray!15},
tick label style={font=\scriptsize},label style={font=\small},title style={font=\small},
ymin=0.71739365,ymax=0.760297207]
\nextgroupplot[title={(a) Accuracy versus tree count},xlabel={Number of trees $T$},
ylabel={Mean test accuracy},xtick={20,40,60,100,200},
legend to name=fairCurveLegend,legend columns=4,legend style={font=\scriptsize,draw=none,
/tikz/every even column/.append style={column sep=3pt}}]
\addplot[color=fairC0,solid,line width=0.8pt] coordinates {(20,0.721545607) (40,0.723651816) (60,0.726046376) (100,0.728771678) (200,0.7302036)};
\addlegendentry{100/0/0\% ($\bar k=1.00$)}
\addplot[color=fairC1,solid,line width=0.8pt] coordinates {(20,0.747903819) (40,0.742645826) (60,0.746474541) (100,0.747768532) (200,0.747574262)};
\addlegendentry{0/100/0\% ($\bar k=2.00$)}
\addplot[color=fairC2,solid,line width=0.8pt] coordinates {(20,0.747405173) (40,0.749170056) (60,0.749182263) (100,0.752152887) (200,0.751334756)};
\addlegendentry{0/0/100\% ($\bar k=3.00$)}
\addplot[color=fairC3,dashed,line width=0.8pt] coordinates {(20,0.727836685) (40,0.730346837) (60,0.731547385) (100,0.732862107) (200,0.734478762)};
\addlegendentry{95/5/0\% ($\bar k=1.05$)}
\addplot[color=fairC4,dashed,line width=0.8pt] coordinates {(20,0.73372712) (40,0.733684112) (60,0.735872107) (100,0.739348705) (200,0.738455662)};
\addlegendentry{90/10/0\% ($\bar k=1.10$)}
\addplot[color=fairC5,dashed,line width=0.8pt] coordinates {(20,0.740635782) (40,0.741210166) (60,0.743275942) (100,0.745725947) (200,0.745488267)};
\addlegendentry{75/25/0\% ($\bar k=1.25$)}
\addplot[color=fairC6,dashed,line width=0.8pt] coordinates {(20,0.745330488) (40,0.744510459) (60,0.746347685) (100,0.746511975) (200,0.74711451)};
\addlegendentry{50/50/0\% ($\bar k=1.50$)}
\addplot[color=fairC7,dotted,line width=0.8pt] coordinates {(20,0.72920301) (40,0.729911507) (60,0.733761174) (100,0.735299154) (200,0.736343946)};
\addlegendentry{95/0/5\% ($\bar k=1.10$)}
\addplot[color=fairC8,dotted,line width=0.8pt] coordinates {(20,0.735337287) (40,0.736434512) (60,0.738485017) (100,0.74149888) (200,0.740500232)};
\addlegendentry{90/0/10\% ($\bar k=1.20$)}
\addplot[color=fairC9,dotted,line width=0.8pt] coordinates {(20,0.744060427) (40,0.7439027) (60,0.749229192) (100,0.75051225) (200,0.751564251)};
\addlegendentry{75/0/25\% ($\bar k=1.50$)}
\addplot[color=fairC10,dotted,line width=0.8pt] coordinates {(20,0.748269825) (40,0.752342111) (60,0.753202605) (100,0.75614525) (200,0.755750931)};
\addlegendentry{50/0/50\% ($\bar k=2.00$)}
\addplot[color=fairC11,dashdotted,line width=0.8pt] coordinates {(20,0.746804221) (40,0.743163816) (60,0.746639317) (100,0.749152299) (200,0.748870817)};
\addlegendentry{0/95/5\% ($\bar k=2.05$)}
\addplot[color=fairC12,dashdotted,line width=0.8pt] coordinates {(20,0.747569298) (40,0.746086328) (60,0.747994449) (100,0.749595989) (200,0.749198589)};
\addlegendentry{0/90/10\% ($\bar k=2.10$)}
\addplot[color=fairC13,dashdotted,line width=0.8pt] coordinates {(20,0.746483461) (40,0.74645805) (60,0.750293248) (100,0.751281199) (200,0.751260101)};
\addlegendentry{0/75/25\% ($\bar k=2.25$)}
\addplot[color=fairC14,dashdotted,line width=0.8pt] coordinates {(20,0.746690633) (40,0.746641789) (60,0.751333666) (100,0.753466409) (200,0.752436573)};
\addlegendentry{0/50/50\% ($\bar k=2.50$)}
\addplot[color=fairC15,solid,line width=0.8pt] coordinates {(20,0.736136319) (40,0.738008505) (60,0.741285904) (100,0.742355309) (200,0.744086104)};
\addlegendentry{85/10/5\% ($\bar k=1.20$)}
\nextgroupplot[title={(b) Accuracy versus fitting work},xmode=log,xlabel={Mean fitting work (s)}]
\addplot[color=fairC0,solid,line width=0.8pt] coordinates {(0.0142270907,0.721545607) (0.0280359555,0.723651816) (0.0413890441,0.726046376) (0.067713161,0.728771678) (0.13382634,0.7302036)};
\addplot[only marks,color=fairC0,mark=*,mark size=1.5pt] coordinates {(0.0142270907,0.721545607)};
\addplot[only marks,color=fairC0,mark=square*,mark size=1.5pt] coordinates {(0.0280359555,0.723651816)};
\addplot[only marks,color=fairC0,mark=triangle*,mark size=1.5pt] coordinates {(0.0413890441,0.726046376)};
\addplot[only marks,color=fairC0,mark=diamond*,mark size=1.5pt] coordinates {(0.067713161,0.728771678)};
\addplot[only marks,color=fairC0,mark=triangle*,mark size=1.5pt,mark options={rotate=180}] coordinates {(0.13382634,0.7302036)};
\addplot[color=fairC1,solid,line width=0.8pt] coordinates {(0.338133776,0.747903819) (0.672754496,0.742645826) (1.00855898,0.746474541) (1.68070227,0.747768532) (3.37126361,0.747574262)};
\addplot[only marks,color=fairC1,mark=*,mark size=1.5pt] coordinates {(0.338133776,0.747903819)};
\addplot[only marks,color=fairC1,mark=square*,mark size=1.5pt] coordinates {(0.672754496,0.742645826)};
\addplot[only marks,color=fairC1,mark=triangle*,mark size=1.5pt] coordinates {(1.00855898,0.746474541)};
\addplot[only marks,color=fairC1,mark=diamond*,mark size=1.5pt] coordinates {(1.68070227,0.747768532)};
\addplot[only marks,color=fairC1,mark=triangle*,mark size=1.5pt,mark options={rotate=180}] coordinates {(3.37126361,0.747574262)};
\addplot[color=fairC2,solid,line width=0.8pt] coordinates {(16.0787606,0.747405173) (32.0231208,0.749170056) (48.0579232,0.749182263) (80.2307394,0.752152887) (160.047967,0.751334756)};
\addplot[only marks,color=fairC2,mark=*,mark size=1.5pt] coordinates {(16.0787606,0.747405173)};
\addplot[only marks,color=fairC2,mark=square*,mark size=1.5pt] coordinates {(32.0231208,0.749170056)};
\addplot[only marks,color=fairC2,mark=triangle*,mark size=1.5pt] coordinates {(48.0579232,0.749182263)};
\addplot[only marks,color=fairC2,mark=diamond*,mark size=1.5pt] coordinates {(80.2307394,0.752152887)};
\addplot[only marks,color=fairC2,mark=triangle*,mark size=1.5pt,mark options={rotate=180}] coordinates {(160.047967,0.751334756)};
\addplot[color=fairC3,dashed,line width=0.8pt] coordinates {(0.0306203155,0.727836685) (0.0605526971,0.730346837) (0.0900091164,0.731547385) (0.148517737,0.732862107) (0.296049975,0.734478762)};
\addplot[only marks,color=fairC3,mark=*,mark size=1.5pt] coordinates {(0.0306203155,0.727836685)};
\addplot[only marks,color=fairC3,mark=square*,mark size=1.5pt] coordinates {(0.0605526971,0.730346837)};
\addplot[only marks,color=fairC3,mark=triangle*,mark size=1.5pt] coordinates {(0.0900091164,0.731547385)};
\addplot[only marks,color=fairC3,mark=diamond*,mark size=1.5pt] coordinates {(0.148517737,0.732862107)};
\addplot[only marks,color=fairC3,mark=triangle*,mark size=1.5pt,mark options={rotate=180}] coordinates {(0.296049975,0.734478762)};
\addplot[color=fairC4,dashed,line width=0.8pt] coordinates {(0.0470039114,0.73372712) (0.0927971859,0.733684112) (0.138473448,0.735872107) (0.229333711,0.739348705) (0.457367181,0.738455662)};
\addplot[only marks,color=fairC4,mark=*,mark size=1.5pt] coordinates {(0.0470039114,0.73372712)};
\addplot[only marks,color=fairC4,mark=square*,mark size=1.5pt] coordinates {(0.0927971859,0.733684112)};
\addplot[only marks,color=fairC4,mark=triangle*,mark size=1.5pt] coordinates {(0.138473448,0.735872107)};
\addplot[only marks,color=fairC4,mark=diamond*,mark size=1.5pt] coordinates {(0.229333711,0.739348705)};
\addplot[only marks,color=fairC4,mark=triangle*,mark size=1.5pt,mark options={rotate=180}] coordinates {(0.457367181,0.738455662)};
\addplot[color=fairC5,dashed,line width=0.8pt] coordinates {(0.0954602178,0.740635782) (0.18909406,0.741210166) (0.283582023,0.743275942) (0.470881029,0.745725947) (0.940747376,0.745488267)};
\addplot[only marks,color=fairC5,mark=*,mark size=1.5pt] coordinates {(0.0954602178,0.740635782)};
\addplot[only marks,color=fairC5,mark=square*,mark size=1.5pt] coordinates {(0.18909406,0.741210166)};
\addplot[only marks,color=fairC5,mark=triangle*,mark size=1.5pt] coordinates {(0.283582023,0.743275942)};
\addplot[only marks,color=fairC5,mark=diamond*,mark size=1.5pt] coordinates {(0.470881029,0.745725947)};
\addplot[only marks,color=fairC5,mark=triangle*,mark size=1.5pt,mark options={rotate=180}] coordinates {(0.940747376,0.745488267)};
\addplot[color=fairC6,dashed,line width=0.8pt] coordinates {(0.176918393,0.745330488) (0.350482679,0.744510459) (0.525248146,0.746347685) (0.874283224,0.746511975) (1.75110832,0.74711451)};
\addplot[only marks,color=fairC6,mark=*,mark size=1.5pt] coordinates {(0.176918393,0.745330488)};
\addplot[only marks,color=fairC6,mark=square*,mark size=1.5pt] coordinates {(0.350482679,0.744510459)};
\addplot[only marks,color=fairC6,mark=triangle*,mark size=1.5pt] coordinates {(0.525248146,0.746347685)};
\addplot[only marks,color=fairC6,mark=diamond*,mark size=1.5pt] coordinates {(0.874283224,0.746511975)};
\addplot[only marks,color=fairC6,mark=triangle*,mark size=1.5pt,mark options={rotate=180}] coordinates {(1.75110832,0.74711451)};
\addplot[color=fairC7,dotted,line width=0.8pt] coordinates {(0.834736148,0.72920301) (1.65236755,0.729911507) (2.4583099,0.733761174) (4.09098873,0.735299154) (8.16412016,0.736343946)};
\addplot[only marks,color=fairC7,mark=*,mark size=1.5pt] coordinates {(0.834736148,0.72920301)};
\addplot[only marks,color=fairC7,mark=square*,mark size=1.5pt] coordinates {(1.65236755,0.729911507)};
\addplot[only marks,color=fairC7,mark=triangle*,mark size=1.5pt] coordinates {(2.4583099,0.733761174)};
\addplot[only marks,color=fairC7,mark=diamond*,mark size=1.5pt] coordinates {(4.09098873,0.735299154)};
\addplot[only marks,color=fairC7,mark=triangle*,mark size=1.5pt,mark options={rotate=180}] coordinates {(8.16412016,0.736343946)};
\addplot[color=fairC8,dotted,line width=0.8pt] coordinates {(1.6538738,0.735337287) (3.2623439,0.736434512) (4.85912989,0.738485017) (8.09331208,0.74149888) (16.0749614,0.740500232)};
\addplot[only marks,color=fairC8,mark=*,mark size=1.5pt] coordinates {(1.6538738,0.735337287)};
\addplot[only marks,color=fairC8,mark=square*,mark size=1.5pt] coordinates {(3.2623439,0.736434512)};
\addplot[only marks,color=fairC8,mark=triangle*,mark size=1.5pt] coordinates {(4.85912989,0.738485017)};
\addplot[only marks,color=fairC8,mark=diamond*,mark size=1.5pt] coordinates {(8.09331208,0.74149888)};
\addplot[only marks,color=fairC8,mark=triangle*,mark size=1.5pt,mark options={rotate=180}] coordinates {(16.0749614,0.740500232)};
\addplot[color=fairC9,dotted,line width=0.8pt] coordinates {(4.05238885,0.744060427) (7.99827919,0.7439027) (12.0238662,0.749229192) (20.0641795,0.75051225) (39.9434322,0.751564251)};
\addplot[only marks,color=fairC9,mark=*,mark size=1.5pt] coordinates {(4.05238885,0.744060427)};
\addplot[only marks,color=fairC9,mark=square*,mark size=1.5pt] coordinates {(7.99827919,0.7439027)};
\addplot[only marks,color=fairC9,mark=triangle*,mark size=1.5pt] coordinates {(12.0238662,0.749229192)};
\addplot[only marks,color=fairC9,mark=diamond*,mark size=1.5pt] coordinates {(20.0641795,0.75051225)};
\addplot[only marks,color=fairC9,mark=triangle*,mark size=1.5pt,mark options={rotate=180}] coordinates {(39.9434322,0.751564251)};
\addplot[color=fairC10,dotted,line width=0.8pt] coordinates {(8.05545022,0.748269825) (16.0055676,0.752342111) (24.0578125,0.753202605) (40.1054761,0.75614525) (79.8751982,0.755750931)};
\addplot[only marks,color=fairC10,mark=*,mark size=1.5pt] coordinates {(8.05545022,0.748269825)};
\addplot[only marks,color=fairC10,mark=square*,mark size=1.5pt] coordinates {(16.0055676,0.752342111)};
\addplot[only marks,color=fairC10,mark=triangle*,mark size=1.5pt] coordinates {(24.0578125,0.753202605)};
\addplot[only marks,color=fairC10,mark=diamond*,mark size=1.5pt] coordinates {(40.1054761,0.75614525)};
\addplot[only marks,color=fairC10,mark=triangle*,mark size=1.5pt,mark options={rotate=180}] coordinates {(79.8751982,0.755750931)};
\addplot[color=fairC11,dashdotted,line width=0.8pt] coordinates {(1.14224961,0.746804221) (2.26456935,0.743163816) (3.37685976,0.746639317) (5.62317326,0.749152299) (11.2393338,0.748870817)};
\addplot[only marks,color=fairC11,mark=*,mark size=1.5pt] coordinates {(1.14224961,0.746804221)};
\addplot[only marks,color=fairC11,mark=square*,mark size=1.5pt] coordinates {(2.26456935,0.743163816)};
\addplot[only marks,color=fairC11,mark=triangle*,mark size=1.5pt] coordinates {(3.37685976,0.746639317)};
\addplot[only marks,color=fairC11,mark=diamond*,mark size=1.5pt] coordinates {(5.62317326,0.749152299)};
\addplot[only marks,color=fairC11,mark=triangle*,mark size=1.5pt,mark options={rotate=180}] coordinates {(11.2393338,0.748870817)};
\addplot[color=fairC12,dashdotted,line width=0.8pt] coordinates {(1.94500366,0.747569298) (3.84230121,0.746086328) (5.72921542,0.747994449) (9.54468063,0.749595989) (18.9888579,0.749198589)};
\addplot[only marks,color=fairC12,mark=*,mark size=1.5pt] coordinates {(1.94500366,0.747569298)};
\addplot[only marks,color=fairC12,mark=square*,mark size=1.5pt] coordinates {(3.84230121,0.746086328)};
\addplot[only marks,color=fairC12,mark=triangle*,mark size=1.5pt] coordinates {(5.72921542,0.747994449)};
\addplot[only marks,color=fairC12,mark=diamond*,mark size=1.5pt] coordinates {(9.54468063,0.749595989)};
\addplot[only marks,color=fairC12,mark=triangle*,mark size=1.5pt,mark options={rotate=180}] coordinates {(18.9888579,0.749198589)};
\addplot[color=fairC13,dashdotted,line width=0.8pt] coordinates {(4.29506241,0.746483461) (8.48193963,0.74645805) (12.7488432,0.750293248) (21.2740008,0.751281199) (42.3739484,0.751260101)};
\addplot[only marks,color=fairC13,mark=*,mark size=1.5pt] coordinates {(4.29506241,0.746483461)};
\addplot[only marks,color=fairC13,mark=square*,mark size=1.5pt] coordinates {(8.48193963,0.74645805)};
\addplot[only marks,color=fairC13,mark=triangle*,mark size=1.5pt] coordinates {(12.7488432,0.750293248)};
\addplot[only marks,color=fairC13,mark=diamond*,mark size=1.5pt] coordinates {(21.2740008,0.751281199)};
\addplot[only marks,color=fairC13,mark=triangle*,mark size=1.5pt,mark options={rotate=180}] coordinates {(42.3739484,0.751260101)};
\addplot[color=fairC14,dashdotted,line width=0.8pt] coordinates {(8.2166656,0.746690633) (16.3278395,0.746641789) (24.5411234,0.751333666) (40.9118952,0.753466409) (81.4953535,0.752436573)};
\addplot[only marks,color=fairC14,mark=*,mark size=1.5pt] coordinates {(8.2166656,0.746690633)};
\addplot[only marks,color=fairC14,mark=square*,mark size=1.5pt] coordinates {(16.3278395,0.746641789)};
\addplot[only marks,color=fairC14,mark=triangle*,mark size=1.5pt] coordinates {(24.5411234,0.751333666)};
\addplot[only marks,color=fairC14,mark=diamond*,mark size=1.5pt] coordinates {(40.9118952,0.753466409)};
\addplot[only marks,color=fairC14,mark=triangle*,mark size=1.5pt,mark options={rotate=180}] coordinates {(81.4953535,0.752436573)};
\addplot[color=fairC15,solid,line width=0.8pt] coordinates {(0.861274561,0.736136319) (1.67963983,0.738008505) (2.54369129,0.741285904) (4.22603385,0.742355309) (8.42151473,0.744086104)};
\addplot[only marks,color=fairC15,mark=*,mark size=1.5pt] coordinates {(0.861274561,0.736136319)};
\addplot[only marks,color=fairC15,mark=square*,mark size=1.5pt] coordinates {(1.67963983,0.738008505)};
\addplot[only marks,color=fairC15,mark=triangle*,mark size=1.5pt] coordinates {(2.54369129,0.741285904)};
\addplot[only marks,color=fairC15,mark=diamond*,mark size=1.5pt] coordinates {(4.22603385,0.742355309)};
\addplot[only marks,color=fairC15,mark=triangle*,mark size=1.5pt,mark options={rotate=180}] coordinates {(8.42151473,0.744086104)};
\addplot[black!55,dashed] coordinates {(0.7,0.71739365) (0.7,0.760297207)};
\end{groupplot}
\end{tikzpicture}
\par\smallskip\scriptsize Composition: \% one-/two-/three-sighted members ($\bar k$: mean horizon).
\par\pgfplotslegendfromname{fairCurveLegend}
% END GENERATED figure2
\caption{Fixed-architecture forests at depth three with all features: every composition spans 20, 40, 60, 100, and 200 members on the complete 57-dataset cohort, averaged over five paired repetitions and then equally across datasets. Circle, square, upward triangle, diamond, and downward triangle mark these sizes in the right panel. The dashed line denotes a 0.7 s mean-fitting-work reference, not a validated per-dataset budget. Composition is stated explicitly because equal average horizons can describe different mixtures. These descriptive curves do not include model-selection cost or show paired confidence intervals.}
\label{fig:tradeoff}
\end{figure*}
% BEGIN FIGURE2_POSTPLOT_TEXT: retain directly after that environment.
Figure~\ref{fig:tradeoff} shows all sixteen depth-three, all-feature
compositions at $T=20,40,60,100,200$. Neither tree count nor the
farther-horizon fraction gives uniformly monotone observed mean
accuracy. Among plotted points with mean fitting work at most $0.7$
s, the pure horizon-two $T=20$ forest has the highest observed mean:
$74.79\%$ at $0.3381$ s. A $T=100$ forest with $25\%$ horizon-two
members has $74.57\%$ accuracy at $0.4709$ s. Sparse mixtures
therefore do not necessarily improve on smaller pure farther-sighted
forests. They offer other tradeoffs: the $T=20$, $25\%$ horizon-two
forest has $74.06\%$ accuracy at $0.0955$ s. Every plotted composition
containing horizon-three members has mean fitting work above $0.7$ s.
This reference is a retrospective comparison of benchmark means, not
a selected budget policy or a per-dataset feasibility guarantee.
Connecting lines do not represent evaluated intermediate models or
certified equal-cost comparisons, and the figure contains no paired
uncertainty. Its fitting-work axis excludes preparation, prediction,
and selection cost.
% END FIGURE2_POSTPLOT_TEXT

\subsection{Training-selected families}
% New figure environment: paths are relative to paper/array_revision/main.tex.
\begin{figure*}[t]
\centering
% BEGIN GENERATED figure3
\definecolor{fairC0}{HTML}{222222}
\definecolor{fairC1}{HTML}{1F77B4}
\definecolor{fairC2}{HTML}{D62728}
\definecolor{fairC3}{HTML}{2B8C55}
\definecolor{fairC4}{HTML}{66A61E}
\definecolor{fairC5}{HTML}{C59A00}
\definecolor{fairC6}{HTML}{8C6D31}
\definecolor{fairC7}{HTML}{7B3294}
\definecolor{fairC8}{HTML}{A64C9C}
\definecolor{fairC9}{HTML}{DC74AB}
\definecolor{fairC10}{HTML}{B35806}
\definecolor{fairC11}{HTML}{008B8B}
\definecolor{fairC12}{HTML}{17A8BD}
\definecolor{fairC13}{HTML}{527A9C}
\definecolor{fairC14}{HTML}{7F7F7F}
\definecolor{fairC15}{HTML}{9D8F11}
\begin{tikzpicture}
\begin{groupplot}[group style={group size=2 by 1,horizontal sep=1.35cm},
width=0.42\textwidth,height=5.2cm,grid=major,grid style={gray!15},
tick label style={font=\scriptsize},label style={font=\small},title style={font=\small}]
\nextgroupplot[title={(a) Accuracy differences},xlabel={Accuracy difference (pp)},
ytick={0,1,2},yticklabels={{Family 1/2 -- RF},{Family 1/2/3 -- RF},{Family 1/2/3 -- 1/2}},y dir=reverse,ymin=-0.4,ymax=2.4]
\addplot+[only marks,mark=*,color=fairC3,error bars/.cd,x dir=both,x explicit] table[x=x,y=y,x error minus=low,x error plus=high] {
x y low high
0.967393297 0 1.03422225 1.33050186
};
\addplot+[only marks,mark=*,color=fairC2,error bars/.cd,x dir=both,x explicit] table[x=x,y=y,x error minus=low,x error plus=high] {
x y low high
1.56929662 1 1.27015996 1.70740143
};
\addplot+[only marks,mark=*,color=fairC7,error bars/.cd,x dir=both,x explicit] table[x=x,y=y,x error minus=low,x error plus=high] {
x y low high
0.60190332 2 0.710278693 1.03119017
};
\addplot[black!40,dashed] coordinates {(0,-0.4) (0,2.4)};
\nextgroupplot[title={(b) Selection and refit costs},ymode=log,ylabel={Mean time or work (s)},
xtick={0,1,2},xticklabels={RF,{Family 1/2},{Family 1/2/3}},
legend to name=fairTunedLegend,legend columns=1,legend style={font=\scriptsize,draw=none}]
\addplot+[only marks,color=fairC1,mark=square*] coordinates {(0,2.27141672) (1,112.079219) (2,2671.04949)};
\addlegendentry{Selection: tree-fit work}
\addplot+[only marks,color=fairC10,mark=*] coordinates {(0,0.0626177376) (1,0.951821055) (2,12.9374083)};
\addlegendentry{Refit: direct wall time}
\addplot+[only marks,color=fairC0,mark=diamond*] coordinates {(0,3.4049916) (1,118.0351) (2,2693.80316)};
\addlegendentry{Sum of measured stage wall times}
\end{groupplot}
\end{tikzpicture}
\par\smallskip\pgfplotslegendfromname{fairTunedLegend}
% END GENERATED figure3
\caption{Training-only selected families on the complete 57-dataset sample and five paired repetitions. Family 1/2 allows horizons one and two; Family 1/2/3 allows all three. Both can select pure forests. Left: mean paired accuracy differences with pointwise, unadjusted $95\%$ dataset-bootstrap intervals, not intervals for the Holm-adjusted signed-rank tests. Right: standalone cross-validation constituent fitting work, direct selected-model refit wall time, and the sum of measured workflow stage wall times. The last measure excludes checkpoint I/O, startup, downloads, and idle gaps; it is not an uninterrupted end-to-end stopwatch. Common validation-bank work is fully charged to every family requiring it.}
\label{fig:tuned}
\end{figure*}

Mean accuracy is $78.99\%$ for the RF family, $79.96\%$ for the
one-/two-sighted family, and $80.56\%$ for the inclusive family
(Figure~\ref{fig:tuned}). Relative to selected RF, the one-/two-sighted
family has a mean gain of $0.97$ points [$-0.07$, $2.30$], median
$0.00617$ points, $29/8/20$ wins/ties/losses, and $p_H=1$.
The inclusive family gains $1.57$ points [$0.30$, $3.28$], with
median zero, $28/11/18$ wins/ties/losses, and $p_H=0.48838$.
Its balanced-accuracy gain is $1.60$ points. Inclusive versus
one-/two-sighted gives $0.60$ points [$-0.11$, $1.63$], median
zero, and $p_H=1$. The positive inclusive mean interval and the
non-rejected signed-rank location null concern different estimands,
not contradictory decisions about one hypothesis. Neither supports
calling the methods equivalent or saying that no mean improvement
was observed.

Selection does not always choose a mixed model. Among the 285
inclusive choices, 133 are pure forests ($61/28/44$ at horizons
one/two/three) and 152 are mixed (146 use two horizons and six use all three).
The one-/two-sighted family selects 168 pure forests and 117
mixtures. Thus the aggregate selected-family gain cannot be
attributed solely to sparse replacement or heterogeneity.
Depth-three/six/unrestricted selection counts are $66/87/132$ for RF,
$90/106/89$ for the one-/two-sighted family, and $104/103/78$ for
the inclusive family. Both all-feature and square-root-feature
settings are selected in every family. The RF label therefore
denotes a library-estimator candidate space including all-feature
bagged CART, not exclusively square-root-feature ensembles.

\subsection{Selection costs and sensitivity of the selected-family effects}
Standalone cross-validation fitting work averages $2.271$, $112.079$,
and $2,671.049$ s for RF, one-/two-sighted, and inclusive families.
These charges include all banks required by each search space,
even when a pure CART forest is ultimately selected. Direct refit
wall times, including construction/bootstrap preparation but not
test prediction, are $0.0626$, $0.9518$, and $12.9374$ s. Their
ratios to RF are $15.20$ and $206.61$ for the two expanded families;
these are refit ratios, not selection-cost ratios. Mean sums of
measured workflow stage wall times are $3.405$, $118.035$, and
$2,693.803$ s. They are not uninterrupted end-to-end times.
Shared banks are fully attributed to each standalone family but
counted only once in the separate physical accounting. Fixed-forest
constituent work and selected-model direct refit wall time have
different boundaries and are not interchangeable.

The inclusive-versus-RF mean effect becomes $1.03$ points
[$-0.27$, $2.90$] with equal weighting of the 48 known families
and $0.53$ points [$-0.91$, $2.77$] on the 38 tasks outside the named
synthetic/constructed set. Both intervals cross zero. The corresponding one-/two-sighted
effects are $0.61$ [$-0.32$, $2.02$] and $0.38$ [$-0.66$, $2.09$]
points. These prespecified analyses show that the estimated tuned
mean benefits are sensitive to benchmark weighting and membership;
they do not establish an absence of benefit or a broadly transferable
accuracy--cost advantage.
% END GENERATED results_text

\section{Discussion}
% BEGIN GENERATED discussion_text
% Replace the entire Discussion body, including text outside old generated markers.
% First complete-results scientific writing pass; independent review still required.

The completed comparison supports conditional predictive benefits rather
than general mixed-forest superiority. Several shallow fixed replacement
contrasts retain small joint-Holm signed-rank values, including both
all-feature replacements and the square-root-feature horizon-three
replacement. Their positive means coexist with zero or small medians
and substantial extra fitting work. Uniformly increasing horizon is
especially unfavorable for the unrestricted all-feature models, whereas
square-root-feature forests show a different pattern. Better local
terminal training impurity is not a certificate of better prediction,
and neither this depth pattern nor a difference between feature policies
identifies a causal mechanism.

Training-only selection changes the practical comparison. The inclusive
family has a positive estimated mean advantage over selected RF, with a
pointwise mean-effect interval above zero. Such a mean gain may be practically
relevant, but its adjusted signed-rank test
does not reject the location null, its median effect is zero, and the
family-weighted intervals and those for the 38 tasks outside the named
synthetic/constructed set cross zero. These summaries answer different questions. They permit a heterogeneous mean
benefit on the full cohort, not a claim of no improvement, equivalence,
or reliable improvement for a typical future task. Because the expanded
families often select pure forests, their aggregate effects evaluate
larger training-selected candidate spaces, not the isolated benefit of
sparse mixtures. Both mixed-capable spaces contain the entire RF space,
but this nesting does not guarantee a held-out improvement. Their larger
selection spaces and search work are part of the comparison.

The shallow experiment remains a capacity-constrained question, not a
recommended universal forest depth. Tree-size regularization studies
motivate conditional benefits of smaller members, particularly in noisy
regression settings \cite{zhou2023trees,duroux2018impact}; they do not
identify depth three as optimal for these bootstrap Gini classifiers.
The fixed CART forests have higher mean accuracy at deeper capacities,
and selected RF chooses depth six or unrestricted depth in most outer
tasks. Thus deeper forests remain relevant and often useful controls.
Growing-tree consistency results do not establish consistency at fixed
depth three \cite{scornet2015consistency}, and shallow boosting results
do not justify an independently fitted bootstrap forest.

Computation remains a substantial limitation. Full standalone selection
work is much larger for the expanded families than for RF, particularly
when horizon-three banks are required, even if the selected refit is
relatively cheap. The observed costs characterize the implementation,
candidate space, and shared hardware load, not every non-greedy method.
Local lookahead is established \cite{murthy1995lookahead,donick2021stepwise};
the additive child recursion gives a worst-case candidate-evaluation
bound exponential in local horizon, rather than the number of complete
split assignments. Final fitting also sums searches over committed
nodes. Fixed horizon is distinct from final depth. Prefix-count reuse
and Cython improve execution without introducing a new optimization
principle or a global certificate. Cached, bounded, and hybrid solvers,
including SPLIT, require objective- and constraint-aligned comparisons
before an accuracy or efficiency ranking is justified
\cite{aglin2020caching,demirovic2022murtree,vanderlinden2023separable,babbar2025nearoptimal}.

The sample is small, previously explored, and includes related and
constructed tasks. Repeated overlapping holdouts are not independent
datasets; known-family weighting addresses only documented dependence.
Complete-case deletion, supplied numeric encodings, rare-class validation
fallbacks, and bounded tuning grids limit the conclusions. Matching common
hyperparameters does not create equal computational budgets, and remaining
threshold, tie, and random-feature-order differences preclude bitwise
algorithmic identity. The descriptive mean-cost curves neither enforce a
per-dataset budget nor account for selection work. No causal diversity
effect, validated resource frontier, large-data scalability, or production
readiness has been established. A prospective resource-constrained study
would need training-only budget selection, final-fit feasibility checks,
complete cost accounting, and larger independent tasks.
% END GENERATED discussion_text

\section{Conclusion}
% BEGIN GENERATED conclusion_text
% Replace the Conclusion body; no historical tuned-versus-fixed result is retained.
Bounded split optimization is established, but its predictive value depends
on final capacity, feature policy, and how models are selected. Sparse
replacements improve specific shallow fixed ensembles at additional cost.
The inclusive training-selected family has a positive benchmark-mean
advantage over selected RF, alongside a zero median, a non-rejected adjusted
signed-rank location null, and sensitivity intervals crossing zero.
Its selections include pure forests, so this is not solely a sparse-mixture
effect. Substantial selection and refit costs prevent interpreting these
results as a general accuracy--cost improvement. Larger independent,
resource-constrained comparisons with aligned modern search methods remain
necessary.
% END GENERATED conclusion_text

\section*{CRediT authorship contribution statement}
Mathias Valla: Conceptualization, Methodology, Software, Investigation, Formal analysis, Writing--original draft, Writing--review and editing.
\section*{Declaration of competing interest}
The author declares no competing interests.
\section*{Acknowledgments}
Work(s) conducted within the Research Chair DIALog under the aegis of the Risk Foundation, a joint initiative by Universit{\'e} Paris-Dauphine PSL and CNP Assurances.
\section*{Data and code availability}
The datasets are publicly available from PMLB \cite{olson2017pmlb}. Code, dataset and environment manifests, repeat-level results, selected hyperparameters, paired analyses, figures, and manuscript sources are available at \url{https://github.com/MathiasValla/non-greedy-decision-trees}. The repository's \texttt{paper/REPRODUCING\_RESULTS.md} documents how to retrieve and reproduce all results in this article. Historical summary tables are separated from the complete new evaluation; historical aggregates without retained raw fits are not used for statistical inference.
\section*{Declaration of generative AI and AI-assisted technologies}
OpenAI Codex (OpenAI) was used to assist with literature discovery, manuscript revision, benchmark and analysis code, and figure-generation code. The Methods describe its use for code and visualization, the unrecorded tool version/model identifier, and the retained software-version records. Numerical figures are computed from retained experimental records using reproducible workflows; no generative image model supplies data or artwork. The author must review the sources, code, results, and text before submission and takes responsibility for the submitted version.

% BEGIN GENERATED bibliography
\begin{thebibliography}{10}
\expandafter\ifx\csname url\endcsname\relax
  \def\url#1{\texttt{#1}}\fi
\expandafter\ifx\csname urlprefix\endcsname\relax\def\urlprefix{URL }\fi
\expandafter\ifx\csname href\endcsname\relax
  \def\href#1#2{#2} \def\path#1{#1}\fi

\bibitem{breiman1984cart}
L.~Breiman, J.~H. Friedman, R.~A. Olshen, C.~J. Stone, Classification and
  Regression Trees, Wadsworth, Belmont, CA, 1984.

\bibitem{zhou2023trees}
S.~Zhou, L.~Mentch,
  \href{https://onlinelibrary.wiley.com/doi/10.1002/sam.11594}{Trees, forests,
  chickens, and eggs: when and why to prune trees in a random forest},
  Statistical Analysis and Data Mining: The ASA Data Science Journal 16~(1)
  (2023) 45--64, first published online 25 August 2022.
\newblock \href {https://doi.org/10.1002/sam.11594}
  {\path{doi:10.1002/sam.11594}}.
\newline\urlprefix\url{https://onlinelibrary.wiley.com/doi/10.1002/sam.11594}

\bibitem{duroux2018impact}
R.~Duroux, E.~Scornet, \href{https://doi.org/10.1051/ps/2018008}{Impact of
  subsampling and tree depth on random forests}, {ESAIM}: Probability and
  Statistics 22 (2018) 96--128.
\newblock \href {https://doi.org/10.1051/ps/2018008}
  {\path{doi:10.1051/ps/2018008}}.
\newline\urlprefix\url{https://doi.org/10.1051/ps/2018008}

\bibitem{quinlan1986induction}
J.~R. Quinlan, Induction of decision trees, Machine Learning 1~(1) (1986)
  81--106.
\newblock \href {https://doi.org/10.1007/BF00116251}
  {\path{doi:10.1007/BF00116251}}.

\bibitem{breiman1996bagging}
L.~Breiman, Bagging predictors, Machine Learning 24~(2) (1996) 123--140.
\newblock \href {https://doi.org/10.1007/BF00058655}
  {\path{doi:10.1007/BF00058655}}.

\bibitem{breiman2001random}
L.~Breiman, Random forests, Machine Learning 45~(1) (2001) 5--32.
\newblock \href {https://doi.org/10.1023/A:1010933404324}
  {\path{doi:10.1023/A:1010933404324}}.

\bibitem{norton1989generating}
S.~W. Norton,
  \href{https://www.ijcai.org/Proceedings/89-1/Papers/128.pdf}{Generating
  better decision trees}, in: Proceedings of the Eleventh International Joint
  Conference on Artificial Intelligence, 1989, pp. 800--805.
\newline\urlprefix\url{https://www.ijcai.org/Proceedings/89-1/Papers/128.pdf}

\bibitem{ragavan1993lookahead}
H.~Ragavan, L.~A. Rendell, Lookahead feature construction for learning hard
  concepts, in: Machine Learning Proceedings 1993, Morgan Kaufmann, 1993, pp.
  252--259.
\newblock \href {https://doi.org/10.1016/B978-1-55860-307-3.50039-3}
  {\path{doi:10.1016/B978-1-55860-307-3.50039-3}}.

\bibitem{murthy1995lookahead}
S.~K. Murthy, S.~L. Salzberg,
  \href{https://www.ijcai.org/Proceedings/95-2/Papers/002.pdf}{Lookahead and
  pathology in decision tree induction}, in: Proceedings of the Fourteenth
  International Joint Conference on Artificial Intelligence, 1995, pp.
  1025--1031.
\newline\urlprefix\url{https://www.ijcai.org/Proceedings/95-2/Papers/002.pdf}

\bibitem{elomaa2005lookahead}
T.~Elomaa, T.~Malinen, On look-ahead and pathology in decision tree learning,
  Journal of Experimental \& Theoretical Artificial Intelligence 17~(1--2)
  (2005) 19--33.
\newblock \href {https://doi.org/10.1080/09528130512331315936}
  {\path{doi:10.1080/09528130512331315936}}.

\bibitem{esmeir2007anytime}
S.~Esmeir, S.~Markovitch,
  \href{https://www.jmlr.org/papers/v8/esmeir07a.html}{Anytime learning of
  decision trees}, Journal of Machine Learning Research 8 (2007) 891--933.
\newline\urlprefix\url{https://www.jmlr.org/papers/v8/esmeir07a.html}

\bibitem{norouzi2015efficient}
M.~Norouzi, M.~D. Collins, M.~A. Johnson, D.~J. Fleet, P.~Kohli, Efficient
  non-greedy optimization of decision trees, in: Advances in Neural Information
  Processing Systems, Vol.~28, 2015, pp. 1729--1737.

\bibitem{donick2021stepwise}
D.~Donick, S.~C. Lera,
  \href{https://www.nature.com/articles/s41598-021-88571-3}{Uncovering feature
  interdependencies in high-noise environments with stepwise lookahead decision
  forests}, Scientific Reports 11 (2021) 9238, article 9238.
\newblock \href {https://doi.org/10.1038/s41598-021-88571-3}
  {\path{doi:10.1038/s41598-021-88571-3}}.
\newline\urlprefix\url{https://www.nature.com/articles/s41598-021-88571-3}

\bibitem{geurts2006extremely}
P.~Geurts, D.~Ernst, L.~Wehenkel,
  \href{https://doi.org/10.1007/s10994-006-6226-1}{Extremely randomized trees},
  Machine Learning 63~(1) (2006) 3--42.
\newblock \href {https://doi.org/10.1007/s10994-006-6226-1}
  {\path{doi:10.1007/s10994-006-6226-1}}.
\newline\urlprefix\url{https://doi.org/10.1007/s10994-006-6226-1}

\bibitem{horn2026randomdepth}
D.~Horn, T.~M. Krabel, T.~N.~T. Tran, A.~Groll, C.~Jentsch,
  \href{https://doi.org/10.1007/s00180-025-01697-0}{The impact of random tree
  depth---a novel randomization process for ensemble methods}, Computational
  Statistics 41 (2026) 25, article 25.
\newblock \href {https://doi.org/10.1007/s00180-025-01697-0}
  {\path{doi:10.1007/s00180-025-01697-0}}.
\newline\urlprefix\url{https://doi.org/10.1007/s00180-025-01697-0}

\bibitem{kuncheva2003diversity}
L.~I. Kuncheva, C.~J. Whitaker,
  \href{https://doi.org/10.1023/A:1022859003006}{Measures of diversity in
  classifier ensembles and their relationship with the ensemble accuracy},
  Machine Learning 51 (2003) 181--207.
\newblock \href {https://doi.org/10.1023/A:1022859003006}
  {\path{doi:10.1023/A:1022859003006}}.
\newline\urlprefix\url{https://doi.org/10.1023/A:1022859003006}

\bibitem{bertsimas2017optimal}
D.~Bertsimas, J.~Dunn, Optimal classification trees, Machine Learning 106~(7)
  (2017) 1039--1082.
\newblock \href {https://doi.org/10.1007/s10994-017-5633-9}
  {\path{doi:10.1007/s10994-017-5633-9}}.

\bibitem{verwer2019learning}
S.~Verwer, Y.~Zhang, Learning optimal classification trees using a binary
  linear program formulation, in: Proceedings of the AAAI Conference on
  Artificial Intelligence, Vol.~33, 2019, pp. 1625--1632.
\newblock \href {https://doi.org/10.1609/aaai.v33i01.33011624}
  {\path{doi:10.1609/aaai.v33i01.33011624}}.

\bibitem{hu2019optimal}
X.~Hu, C.~Rudin, M.~Seltzer, Optimal sparse decision trees, in: Advances in
  Neural Information Processing Systems, Vol.~32, 2019, pp. 7265--7273.

\bibitem{lin2020generalized}
J.~Lin, C.~Zhong, D.~Hu, C.~Rudin, M.~Seltzer,
  \href{https://proceedings.mlr.press/v119/lin20g.html}{Generalized and
  scalable optimal sparse decision trees}, in: Proceedings of the 37th
  International Conference on Machine Learning, Vol. 119 of Proceedings of
  Machine Learning Research, 2020, pp. 6150--6160.
\newline\urlprefix\url{https://proceedings.mlr.press/v119/lin20g.html}

\bibitem{aglin2020caching}
G.~Aglin, S.~Nijssen, P.~Schaus,
  \href{https://ojs.aaai.org/index.php/AAAI/article/view/5711}{Learning optimal
  decision trees using caching branch-and-bound search}, in: Proceedings of the
  AAAI Conference on Artificial Intelligence, Vol.~34, 2020, pp. 3146--3153.
\newblock \href {https://doi.org/10.1609/aaai.v34i04.5711}
  {\path{doi:10.1609/aaai.v34i04.5711}}.
\newline\urlprefix\url{https://ojs.aaai.org/index.php/AAAI/article/view/5711}

\bibitem{demirovic2022murtree}
E.~Demirovi{\'c}, A.~Lukina, E.~Hebrard, J.~Chan, J.~Bailey, C.~Leckie,
  K.~Ramamohanarao, P.~J. Stuckey,
  \href{https://jmlr.org/papers/v23/20-520.html}{{MurTree}: Optimal decision
  trees via dynamic programming and search}, Journal of Machine Learning
  Research 23~(26) (2022) 1--47.
\newline\urlprefix\url{https://jmlr.org/papers/v23/20-520.html}

\bibitem{vanderlinden2023separable}
J.~G.~M. van~der Linden, M.~M. de~Weerdt, E.~Demirovi{\'c},
  \href{https://proceedings.neurips.cc/paper_files/paper/2023/hash/1d5fce9627e15c84db572a66e029b1fc-Abstract-Conference.html}{Necessary
  and sufficient conditions for optimal decision trees using dynamic
  programming}, in: Advances in Neural Information Processing Systems, Vol.~36,
  2023, pp. 9173--9212.
\newblock \href {https://doi.org/10.52202/075280-0404}
  {\path{doi:10.52202/075280-0404}}.
\newline\urlprefix\url{https://proceedings.neurips.cc/paper_files/paper/2023/hash/1d5fce9627e15c84db572a66e029b1fc-Abstract-Conference.html}

\bibitem{aghaei2025strong}
S.~Aghaei, A.~G{\'o}mez, P.~Vayanos, Strong optimal classification trees,
  Operations Research 73~(4) (2025) 2223--2241.
\newblock \href {https://doi.org/10.1287/opre.2021.0034}
  {\path{doi:10.1287/opre.2021.0034}}.

\bibitem{mctavish2022reference}
H.~McTavish, C.~Zhong, R.~Achermann, I.~Karimalis, J.~Chen, C.~Rudin,
  M.~Seltzer,
  \href{https://ojs.aaai.org/index.php/AAAI/article/view/21194}{Fast sparse
  decision tree optimization via reference ensembles}, in: Proceedings of the
  AAAI Conference on Artificial Intelligence, Vol.~36, 2022, pp. 9604--9613.
\newblock \href {https://doi.org/10.1609/aaai.v36i9.21194}
  {\path{doi:10.1609/aaai.v36i9.21194}}.
\newline\urlprefix\url{https://ojs.aaai.org/index.php/AAAI/article/view/21194}

\bibitem{babbar2025nearoptimal}
V.~Babbar, H.~McTavish, C.~Rudin, M.~Seltzer,
  \href{https://proceedings.mlr.press/v267/babbar25a.html}{Near-optimal
  decision trees in a {SPLIT} second}, in: Proceedings of the 42nd
  International Conference on Machine Learning, Vol. 267 of Proceedings of
  Machine Learning Research, PMLR, 2025, pp. 2114--2175.
\newline\urlprefix\url{https://proceedings.mlr.press/v267/babbar25a.html}

\bibitem{vanderlinden2026optimalgreedy}
J.~G.~M. van~der Linden, D.~Vos, M.~M. de~Weerdt, S.~Verwer, E.~Demirovi{\'c},
  \href{https://openreview.net/forum?id=DvDOAtskXl}{Optimal or greedy decision
  trees? revisiting their objectives, tuning, and performance}, Transactions on
  Machine Learning Research (2026).
\newline\urlprefix\url{https://openreview.net/forum?id=DvDOAtskXl}

\bibitem{scornet2015consistency}
E.~Scornet, G.~Biau, J.-P. Vert,
  \href{https://doi.org/10.1214/15-AOS1321}{Consistency of random forests}, The
  Annals of Statistics 43~(4) (2015) 1716--1741.
\newblock \href {https://doi.org/10.1214/15-AOS1321}
  {\path{doi:10.1214/15-AOS1321}}.
\newline\urlprefix\url{https://doi.org/10.1214/15-AOS1321}

\bibitem{olson2017pmlb}
R.~S. Olson, W.~La~Cava, P.~Orzechowski, R.~J. Urbanowicz, J.~H. Moore, {PMLB}:
  a large benchmark suite for machine learning evaluation and comparison,
  BioData Mining 10 (2017) 36.
\newblock \href {https://doi.org/10.1186/s13040-017-0154-4}
  {\path{doi:10.1186/s13040-017-0154-4}}.

\bibitem{cawley2010selection}
G.~C. Cawley, N.~L.~C. Talbot,
  \href{https://jmlr.org/papers/v11/cawley10a.html}{On over-fitting in model
  selection and subsequent selection bias in performance evaluation}, Journal
  of Machine Learning Research 11 (2010) 2079--2107.
\newline\urlprefix\url{https://jmlr.org/papers/v11/cawley10a.html}

\bibitem{demsar2006comparisons}
J.~Dem{\v s}ar, \href{https://jmlr.org/papers/v7/demsar06a.html}{Statistical
  comparisons of classifiers over multiple data sets}, Journal of Machine
  Learning Research 7 (2006) 1--30.
\newline\urlprefix\url{https://jmlr.org/papers/v7/demsar06a.html}

\bibitem{holm1979sequential}
S.~Holm, \href{https://www.jstor.org/stable/4615733}{A simple sequentially
  rejective multiple test procedure}, Scandinavian Journal of Statistics 6~(2)
  (1979) 65--70.
\newline\urlprefix\url{https://www.jstor.org/stable/4615733}

\end{thebibliography}
% END GENERATED bibliography
\end{document}
